
This section is based on the work on \cite{majumdar2007randommatricesulamproblem}.

A very commom - and wildely important - method used in evolutionary biology, and man correlated areas, is the alignement of sequences. The goal of this method is to search for similarities in patters in different sequences. A classic, and commonly studied alignment problem is the so called 'longest common subsequence' (LCS). The setup is of two n- sequence and $m$- seuquence, composed of $n$ and $m$ letters, respectively, of the same $c$-alphabet (a set of $c$ letters).  This set could be, for example, the set of basis of the DNA.  Consider the sequences $\alpha, \beta$ given by
\begin{align*}
	\alpha \quad & : \qquad A \quad C \quad G \quad C \quad T \quad A \quad C \\
	\beta \quad & : \qquad C \quad T \quad G \quad A \quad C
\end{align*}
A subsequence of $\alpha$ is an ordered sublist of $\alpha$ (entries of which need not be consecutive in $\alpha$), e.g, $\{C, G, C, T\}$, but not $\{T, G, C\}$. A common subsequence of two sequences $\alpha$ and $\beta$ is a subsequence of both of them. For example, the subsequence $\{C, G, A, C\}$ is a common subsequence of both $\alpha$ and $\beta$. The aim of the LCS problem is to find the longest of such common subsequences between two fixed sequences $\alpha$ and $\beta$. 

A particularly important application of the LCS problem is to quantify the closeness between two DNA sequences. In evolutionary biology, the genes responsible for building specifci proteins evolve with time and by fniding the LCS of the same gene in different species, one can learn what has been conserved in time. Also, when a new DNA molecule is sequenced in vitro, it is important to know whether it is really new or it already exists. This is achieved quantitatively by measuring the LCS of the new molecule with another existing already in the database. 

For a pair of fixed sequences of length $i$ and $j$ respectively, the length $L_{i,j}$ of their LCS is just a number. However, in the stochastic version of the LCS problem one compares two random sequences drawn from $c$-alphabets and hence the length $L_{i,j}$ is a random variable. A major challenge over the last three decades has been to determine the statistics of $L_{i,j}$. For equally long sequences $(i = j = n)$, it has been proved that $\langle L_{n,n} \rangle \approx \gamma c n$ for $n \gg 1$, where the averaging is performed over all realizations of the random sequences. The constant $\gamma_c$ is known as the Chvátal-Sankoff constant which, to date, remains undetermined though there exists several bounds, a conjecture due to Steele that $\gamma_c = 2 / (1 + \sqrt{n})$ and a recent proof that $\gamma_c \rightarrow 2 / \sqrt{c} $ as $c \rightarrow \infty$. Unfortunately, no exact results are available for the finite size corrections to the leading behavior of the average $\langle L_{n,n} \rangle$, for the variance, and also for the full probability distribution of $L{n,n}$. Thus, despite tremendous analytical and numerical efforts, exact solution of the random LCS problem has, so far, remained elusive. Therefore it is important to fnid other variants of this LCS problem that may be analytically tractable. 

Computationally, the easiest way to determine the length $L_{i,j}$ of the LCS of two arbitrary sequences of lengths $i$ and $j$ (in polynomial time $O(ij)$) is via using the recursive algorithm
\begin{equation}
	L_{i,j} = \max[L_{i-1,j}, L_{i,j-1}, L_{i-1,j-1}+ \eta_{i,j}] 
\end{equation}
subject to the initial conditions $L_{0,j} = L_{i,0} = L_{0,0} = 0$. The variable $\eta_{i,j}$ is either $1$ when the characters at the positions $i$ (in the sequence $\alpha$) and $j$ (in the sequence $\beta$) match each other, or $0$ if they do not. Note that the variables $\eta_{i,j}$’s are not independent of each other. These residual correlations between the $\eta_{i,j}$ variables make the LCS problem rather complicated. Note however that for two random sequences drawn from $c$ alphabets, these correlations between the $\eta_{i,j}$ variables vanish in the $c \rightarrow \infty$ limit. 

A natural question is how important are these correlations between the $\eta_{i,j}$ variables, e.g., do they affect the asymptotic statistics of $L_{i,j}$’s for large $i$ and $j$? Is the problem solvable if one ignores these correlations? These questions naturally lead to the Bernoulli matching (BM) model which is a simpler variant of the original LCS problem where one ignores the correlations between $\eta_{i,j}$’s for all $c$. The length LBM $i,j$ of the BM model satisfeis the same recursion relation as before, except that $\eta_{i,j}$’s are now independent and each drawn from the bimodal distribution: $\p(\eta) = (1/c) \delta_{\eta,1} + (1-1/c) \delta_{\eta,0}$. This approximation is expected to be exact only in the appropriately taken $c \rightarrow \infty$ limit. Nevertheless, for finite $c$, the results on the BM model can serve as a useful benchmark for the original LCS model to decide if indeed the correlations between $\eta_{i,j}$’s are important or not. Unfortunately, even in the absence of correlations, the exact aymptotic distribution of $LBM_{i,j}$ in the BM model has so far remained elusive, mainly due to the nonlinear nature of the recursion relation. 

The reference presents present an exact asymptotic formula for the distribution of the length $LBM_{n,n}$ in the BM model for all $c$. So far, only the leading asymptotic behavior of the average length in the BM model was known using the ‘cavity’ method of spin glass physics,
\begin{equation}
	\langle L^{(BM)}_{n,n} \rangle \approx \gamma^{(BM)}_c n
\end{equation}
where $\gamma_c^{(BM)} = 2/(1+\sqrt{c})$, the same as the conjectured value of the Chvátal-Sankoff constant $\gamma_c$ for the original LCS model. However, other properties such as the variance or the distribution of $LBM_{n,n}$ remained untractable even in the BM model. We have shown, as illustrated below, that for large $n$,
\begin{equation}
	LBM_{n,n} \rightarrow \gamma^{(BM)}_c n + f(c) n^{1/3} \chi	
\end{equation}
where $\chi$ is a random variable with a $n$-independent distribution, $\p(\chi \leq x) = F_2(x)$, which is precisely the Tracy-Widom distribution.

Thus, the Tracy-Widom distribution also describes the asymptotic distribution of the optimal matching length in the BM model, for all $c$. Given that the correlations in the original LCS model become negligible in the $c \rightarrow \infty$ limit, it is likely that the BM asymptotics would also hold for the original LCS model in the $c \rightarrow \infty$. An important open problem is to determine whether the Tracy-Widom distribution also appears in the LCS problem for finite $c$. The distribution obtained for all $c$ in the BM model can serve as a useful benchmark to which future simulations of the LCS problem can be compared. 